\section{Explicit descent of the multiplier}
\label{sec:descent}

Fix a prime $p$. A valid state is a pair $s=(m,x)$ satisfying
\begin{equation}
\label{eq:valid}
 I_p(s):\quad 1\le m<p\quad\land\quad \Nrm{x}=mp.
\end{equation}
The task is to reach $m=1$ while preserving this assertion. Instead of
choosing a least possible multiplier, we define the next state directly.

\subsection{Even multipliers: exact parity halving}

\begin{lemma}[Halving an even norm]
\label{lem:half}
There is a primitive recursive function $H:\ZZ^4\to\ZZ^4$ such that
\[
 2\mid\Nrm{x}\quad\Longrightarrow\quad 2\Nrm{H(x)}=\Nrm{x}.
\]
\end{lemma}

\begin{proof}
Since $z^2\equiv z\pmod2$, an even norm means that the number of odd
coordinates is even. For $x=(a,b,c,d)$ choose one of three orders:
\[
 (a,b,c,d)\quad\hbox{if }a\equiv b\pmod2;
\]
otherwise choose $(a,c,b,d)$ if $a\equiv c\pmod2$, and $(a,d,b,c)$ in the
remaining case. These are explicit parity tests.

In the first two cases, one chosen pair has equal parity, and the parity of
the total norm forces the remaining pair to have equal parity as well.
In the last case $b$ and $c$ both have parity opposite to $a$, so they
have equal parity. The parity of the total then forces $d$ to have the
same parity as $a$. Hence each chosen pair has equal parity in all cases.

Write the reordered tuple as $(u,v,w,z)$ and define
\begin{equation}
\label{eq:half}
 H(x)=\left(\frac{u+v}{2},\frac{u-v}{2},
              \frac{w+z}{2},\frac{w-z}{2}\right).
\end{equation}
All four divisions are exact, including those of negative numerators.
For either pair,
\[
 (u+v)^2+(u-v)^2=2u^2+2v^2.
\]
Substitute the exact quotient identities, add the two pair identities,
and cancel the positive integer $2$. The result is
$2\Nrm{H(x)}=u^2+v^2+w^2+z^2=\Nrm{x}$.
\end{proof}

If a valid state has even $m>1$, set
\[
 m'=m/2,\qquad x'=H(x).
\]
Then $1\le m'<m$, and Lemma~\ref{lem:half} gives $\Nrm{x'}=m'p$.
Thus this branch strictly decreases the multiplier.

\subsection{Odd multipliers: centered residues and exact division}

Assume now that $m>1$ is odd. For an integer $z$, let $t=z\bmod m$,
where $0\le t<m$, and define
\begin{equation}
\label{eq:center}
 \rho_m(z)=
 \begin{cases}
 t,&2t\le m,\\
 t-m,&2t>m.
 \end{cases}
\end{equation}
Because $m$ is odd, this gives
\begin{equation}
\label{eq:centerbound}
 \rho_m(z)\equiv z\pmod m,\qquad
 \abs{\rho_m(z)}\le\frac{m-1}{2}.
\end{equation}
Put $r_i=\rho_m(x_i)$ and $r=(r_0,r_1,r_2,r_3)$.

\begin{lemma}[The new multiplier]
\label{lem:newmult}
For a valid state with odd $m>1$, the quotient
\[
 k=\Nrm{r}\mathbin{\mathrm{div}}m
\]
is exact and satisfies $1\le k<m$.
\end{lemma}

\begin{proof}
Congruence of the coordinates implies congruence of their squares and
their sums. Therefore
\[
 \Nrm{r}\equiv\Nrm{x}=mp\equiv0\pmod m,
\]
so $\Nrm{r}=mk$. Moreover,
\[
 mk=\Nrm{r}
 \le4\left(\frac{m-1}{2}\right)^2
 =(m-1)^2<m^2.
\]
Cancellation of $m>0$ gives $k<m$. The displayed estimate can be written
without fractions by setting $m=2h+1$: it is simply
$\Nrm{r}\le4h^2<(2h+1)^2$.

It remains to rule out zero. If $k=0$, the nonnegative squares in
$\Nrm{r}=0$ must all be zero. Thus all $r_i=0$, and $m$ divides every
$x_i$. Compute the exact quotients $w_i=x_i/m$. Then
\[
 m^2\Nrm{w}=\Nrm{x}=mp,
\]
and cancellation gives $m\Nrm{w}=p$. This would make $m$ a divisor of
$p$ with $1<m<p$, contradicting the primality test. Since the natural
number $k$ is decidably either zero or positive, it follows that $k\ge1$.
\end{proof}

\begin{lemma}[The new representing tuple]
\label{lem:quotient}
Under the assumptions of Lemma~\ref{lem:newmult}, every coordinate of
$z=x*\bar r$ is divisible by $m$, and the exact quotient tuple
$x'=z/m$ satisfies $\Nrm{x'}=kp$.
\end{lemma}

\begin{proof}
For clarity, with $x=(a,b,c,d)$ and $r=(r,s,t,u)$, the four numerators are
\begin{align*}
 z_0&= ar+bs+ct+du,\\
 z_1&=-as+br-cu+dt,\\
 z_2&=-at+bu+cr-ds,\\
 z_3&=-au-bt+cs+dr.
\end{align*}
Modulo $m$ one can replace $a,b,c,d$ by $r,s,t,u$, respectively.
The first expression becomes $r^2+s^2+t^2+u^2=\Nrm{r}=mk$;
each of the other three becomes zero by cancellation of paired products.
Thus $m\mid z_i$ for each $i$, with exact signed quotient $x'_i$.

Euler's identity and invariance of the norm under conjugation give
\[
 m^2\Nrm{x'}=\Nrm{z}
 =\Nrm{x}\Nrm{\bar r}
 =(mp)(mk)=m^2pk.
\]
Cancel the positive integer $m^2$ to obtain $\Nrm{x'}=pk$.
\end{proof}

\subsection{The total transition and its decrease theorem}

The transition on coded states is now completely specified:
\begin{equation}
\label{eq:step}
 \Step(m,x)=
 \begin{cases}
 (m,x),&m\le1,\\
 (m/2,H(x)),&m>1\text{ even},\\
 (\Nrm{r}/m,(x*\bar r)/m),&m>1\text{ odd},
 \end{cases}
\end{equation}
where $r$ is given by~\eqref{eq:center}. Quotient operations are defined
on all inputs; their exactness is a proved property of valid states.
Therefore this is a total primitive recursive function even on malformed
state codes. It does not need a partial ``divide if divisible'' operation.

\begin{proposition}[One-step descent]
\label{prop:step}
If $p$ is prime and $I_p(m,x)$ holds, then $I_p(\Step(m,x))$ holds.
If $m>1$, the new multiplier $m'$ satisfies $1\le m'<m$.
If $m=1$, the state is unchanged.
\end{proposition}

\begin{proof}
The terminal case follows from the first branch of~\eqref{eq:step}.
The even case follows from Lemma~\ref{lem:half}, with $m'=m/2$.
The odd case follows from Lemmas~\ref{lem:newmult} and~\ref{lem:quotient}.
These exhaustive cases are determined by natural-number comparison and
parity, both decidable arithmetic tests.
\end{proof}

\begin{remark}[Why the parity branch matters]
A centered-residue bound for an arbitrary even modulus gives only
$\abs{r_i}\le m/2$, hence $\Nrm{r}/m\le m$. That estimate alone does
not prove strict descent. Treating even multipliers by exact halving removes
the boundary case altogether. For the remaining odd multipliers,
$\abs{r_i}\le(m-1)/2$ yields strict decrease directly. The separate
argument excluding $k=0$ is also essential.
\end{remark}

\section{Termination by a fixed primitive recursive clock}
\label{sec:clock}

We record the precise induction used instead of an appeal to infinite
descent.

\begin{lemma}[An absorbing bounded descent]
\label{lem:clock}
Let $D$ be a primitive recursive operation on natural state codes, $J$ a
quantifier-free state predicate, and $R$ a natural-valued primitive
recursive rank. Suppose arithmetic proves
\[
 J(s)\to J(D(s)),\qquad
 J(s)\to R(D(s))\le R(s)\monus1.
\]
Define $D^0(s)=s$ and $D^{t+1}(s)=D(D^t(s))$ by primitive recursion.
Then quantifier-free induction proves
\begin{equation}
\label{eq:rankinvariant}
 J(s)\to
 \bigl(J(D^t(s))\land R(D^t(s))\le R(s)\monus t\bigr).
\end{equation}
In particular, $R(D^t(s))=0$ whenever $t\ge R(s)$ and $J(s)$ holds.
\end{lemma}

\begin{proof}
For $t=0$, both claims follow from $D^0(s)=s$ and $R(s)\monus0=R(s)$.
Assume the displayed invariant at $t$. Preservation gives $J(D^{t+1}(s))$,
and the rank inequality gives
\[
 \begin{split}
 R(D^{t+1}(s))
 &\le R(D^t(s))\monus1\\
 &\le(R(s)\monus t)\monus1
   =R(s)\monus(t+1).
 \end{split}
\]
Monotonicity and the last identity are elementary laws of truncated
subtraction, proved by quantifier-free induction. Formula
\eqref{eq:rankinvariant}, with state code $s$ as a parameter, is itself
quantifier-free. At $t\ge R(s)$ its right-hand bound is zero.
\end{proof}

For the prime descent, use
\[
 J(s)=I_p(s),\qquad R(m,x)=m\monus1.
\]
If $m=1$, the rank is zero and remains zero. If $m>1$, strict multiplier
decrease gives the required saturated decrease of the rank. Define
\begin{equation}
\label{eq:primeprog}
 s_t(p)=\Step^{\,t}(m_0(p),x^{(0)}(p)),\qquad
 P(p)=\text{the tuple component of }s_p(p).
\end{equation}
Here the clock is the given input $p$, not the first time a test succeeds.

\begin{theorem}[Prime witness]
\label{thm:prime}
The function $P$ is primitive recursive, and $\ipra$ proves
\[
 \Prime(p)\to\Nrm{P(p)}=p.
\]
\end{theorem}

\begin{proof}
The seed is valid by Proposition~\ref{prop:seed}. Apply
Lemma~\ref{lem:clock} and Proposition~\ref{prop:step} to obtain
\[
 I_p(s_t(p)),\qquad
 (m_t\monus1)\le(m_0\monus1)\monus t.
\]
These assertions are obtained by quantifier-free induction, with the
prime-test equation as a parameter assumption. Since $m_0<p$, at $t=p$
the rank is zero. Validity also gives $m_p\ge1$, so $m_p=1$. The norm
part of validity is now $\Nrm{P(p)}=p$.

The seed, transition, and iteration have already been defined as total
primitive recursive functions. The proof does not establish their
termination by a separate unbounded search; it establishes correctness
after a predetermined number of steps.
\end{proof}

\subsection{Two fully worked descent traces}

For $p=7$, the lexicographic seed search returns $(a,b)=(2,3)$, so
\[
 2\cdot7=2^2+3^2+1^2+0^2.
\]
The tuple is reordered as $(2,0,3,1)$ into equal-parity pairs. Halving
sums and differences gives $(1,1,2,1)$ and
$7=1^2+1^2+2^2+1^2$.

For $p=23$, the seed is $(2,8,1,0)$ with multiplier $3$. Centering modulo
$3$ gives $r=(-1,-1,1,0)$ and $\Nrm{r}=3$. Hence $k=1$.
The conjugate product is
\[
 (2,8,1,0)*(-1,1,-1,0)=(-9,-6,-3,-9).
\]
Divide by $3$ to obtain $(-3,-2,-1,-3)$, of norm $23$.

A trace exercising both branches is
\[
 \begin{array}{c|c|c}
 p=31&m&x\\\hline
 \text{seed}&6&(4,13,1,0)\\
 \text{even halving}&3&(2,2,7,6)\\
 \text{odd descent}&1&(1,2,-1,-5).
 \end{array}
\]
Every row satisfies $\Nrm{x}=31m$.
